\documentclass{jsarticle}
\usepackage[dvipdfmx]{graphicx}
\usepackage{ascmac} % itemboxを利用するために必要
\usepackage{url}

%図にEPSファイルを採用した場合，PDFファイルが生成されないことがあるので，PNGやJPEG，PDFファイルを図に採用したほうがよい．

\title{ソフトウェア設計及び実験\\前期・個人開発進捗報告書}

\author{学生番号：6125010806　氏名：清水 統貴}
\date{2026年7月8日}

\begin{document}
\maketitle

\section{開発期間2025年7月2日～7月8日の進捗状況}
%

\subsection{前回からの課題と進展事項}
%
現在までにプレイヤーキャラクターのアクションを完成させ，ステージに背景を追加した．
また，プレイヤーが右に進むと，カメラもスクロールするようにした。ステージ上にギミックを配置
し，少しずつステージづくりに取り掛かっている．具体的な内容は表\ref{table1}に示す．

\begin{table}[htb]
\begin{center}
\caption{課題と進展}
\label{table1}
\begin{tabular}{|p{7cm}|p{7cm}|}\hline
課題 & 進展 \\ \hline \hline
ブースト機能の追加 & ゲージを消費しながらプレイヤーが縦横無尽にステージを飛び回れるアクションを追加した． \\ \hline
キーボード操作を実装 & テストをしやすくするためにこの段階でキーボード操作を実装した． \\ \hline
横スクロール機能の実装 & プレイヤーが右に進むとそれにカメラがついてくるようにした．また，スーパーマリオブラザーズ同様，右に進むと画面外になった左には戻れないようにした． \\ \hline
ギミックの配列管理とステージ配置 & Gimmick構造体を用意し，ギミックを配列で管理するシステムを構築．また，指定した座標にギミックを配置する関数を追加． \\ \hline
当たり判定の実装 & プレイヤーとギミックの重なりを検知する関数を実装．また，ギミックに応じた効果を受ける関数を実装した． \\ \hline
ダメージを受けた際の無敵タイマーとノックバックの実装 & ダメージを受けたときに一時的に無敵になる無敵タイマーを実装した．また，トゲに接触した際，プレイヤーを弾き飛ばすノックバックを接触方向に応じて発生させる処理を追加した． \\ \hline 
\end{tabular}
\end{center}
\end{table}

\begin{figure}[htb]
\centering
\label{figure1}
\includegraphics[scale=0.3]{window.png}
\caption{ゲームを実行したときの様子}
\end{figure}

\section{発生した課題と解決策}
今回の実装過程において，不具合及びエラーが発生したが，原因を特定しすべて解決した．
不具合と解決策は表\ref{table2}を参照．
\begin{table}[htb]
\begin{center}
\caption{課題と進展}
\label{table2}
\begin{tabular}{|p{7cm}|p{7cm}|}\hline
不具合 & 解決策 \\ \hline \hline
ヘッダーファイルの多重インクルードエラー & 各ヘッダーにインクルードガードを追記 \\ \hline
多段ヒットによる即死 & ノックバックと無敵時間を追加することで多段ヒットを防止した． \\ \hline
ノックバックによるカメラ左端の見えない壁のすり抜け & カメラの座標補正がギミックのノックバック処理より前におこなわれていたため，一番最後に行うようにした． \\ \hline
ゲーム実行時のコアダンプ & 終了処理の順番が生成時と逆になっていたためメモリ破壊がおこっていた．順序を正しい順番に修正した． \\ \hline
ダメージを受けた際の無敵タイマーとノックバックの実装 & ダメージを受けたときに一時的に無敵になる無敵タイマーを実装した．また，トゲに接触した際，プレイヤーを弾き飛ばすノックバックを接触方向に応じて発生させる処理を追加した． \\ \hline 
\end{tabular}
\end{center}
\end{table}


\section{開発期間2026年7月9日～15日の予定}

\subsection{完了させる事項}
ステージに配置するギミックを考えて作成していきたい．

\subsection{新しく取り組む事項}
プレイヤーやギミックのテクスチャをきちんとした画像に変更したい．タイトル画面やゴール後の処理についても
実装していきたい．


\begin{thebibliography}{9}
\bibitem{gemini} Google : Gemini \url{https://gemini.google.com/?hl=ja} (2026年7月8日閲覧)
\end{thebibliography}
\end{document}

